% !TEX root = ../main.tex
\chapter{Entorn de desenvolupament i infraestructura de còmput}
\label{chap:infraestructura}

El Capítol~\ref{chap:metodologia} exigeix centenars de cel·les independents, motors V8 reciclables i execucions de dies. Aquest capítol documenta l'entorn que fa viable aquella execució ---maquinari, xarxa, monitoratge, dependències i organització--- i sobre el qual el Capítol~\ref{chap:implementacio} implementa coordinador, agents i connector amb el \emph{ladder}. El coordinador mateix es descriu allà; aquí es fixa on i amb quins recursos corre.

Cal distingir dos volums. La \textbf{matriu dirigida} de 28 agents suma \textbf{6,4 milions} de batalles (6.409.000). El \textbf{còmput total} ---proves, altres formats, entrenament PPO, repeticions i execucions de diagnòstic--- s'aproxima a \textbf{52,7 milions} de batalles. Les dues xifres no són intercanviables: la primera és l'experiment; la segona, el cost de còmput acumulat.

\section{Organització del projecte i gestió amb Kanban}
\label{sec:kanban}

El desenvolupament del projecte s'ha organitzat mitjançant una metodologia àgil basada en un tauler \textbf{Kanban}, dividit en sis seccions funcionals per controlar l'abast i prioritzar el treball:

\begin{itemize}
    \item \textbf{\textit{Aim of the study}:} Defineix els objectius principals, hipòtesis de recerca i el disseny dels 5 paradigmes.
    \item \textbf{\textit{Tentative (If we have time)}:} Línies d'investigació secundàries o exploratòries (com els agents LLM i proves multigeneracionals).
    \item \textbf{\textit{To Do (Icebox/later)}:} Idees i tasques congelades a llarg termini per evitar la dispersió de l'abast (\textit{scope creep}).
    \item \textbf{\textit{To Do}:} Tasques pendents immediates i prioritzades de codi, simulació o redacció.
    \item \textbf{\textit{In Progress}:} Treball actiu en desenvolupament, entrenament PPO o execució de la matriu.
    \item \textbf{\textit{Done}:} Mòduls finalitzats, simulacions executades, dades validades i seccions redactades.
\end{itemize}

Aquesta estructura ha permès controlar l'abast (heurístiques, cerca, imitació, PPO, LLM i redacció) i traçar-lo amb el repositori (Secció~\ref{sec:git_metrics}).

\section{Evolució de l'arquitectura}

\subsection{Fase 1: prototipatge en portàtil}

El desenvolupament inicial, la construcció de les primeres heurístiques (\texttt{v1}--\texttt{v6}) i la redacció de la memòria es van dur a terme en un portàtil:

\begin{itemize}
    \item \textbf{Processador (CPU):} Intel Core Ultra 7 155H (16 nuclis / 22 fils).
    \item \textbf{Memòria RAM:} 16~GB ---suficient per a prototipatge, insuficient per a múltiples instàncies de Showdown en paral·lel.
    \item \textbf{Emmagatzematge:} SSD NVMe d'1~TB per a \path{data/}.
\end{itemize}

Aquesta fase va identificar el coll d'ampolla de RAM i va motivar les neteges amb \texttt{gc.collect()}, els lots continguts i el reinici periòdic de Node.js ja fixats com a principi al Capítol~\ref{chap:metodologia}.

\subsection{Fase 2: estació de treball}

Els experiments finals i la matriu s'executen en una torre configurada per a càrrega sostinguda:

\begin{itemize}
    \item \textbf{Processador (CPU):} AMD Ryzen 7 5700X3D (8 nuclis / 16 fils, 96~MB L3 3D V-Cache), adequat a simulació de partides i a arbres Minimax/MCTS.
    \item \textbf{Refrigeració:} dissipador d'aire Noctua; la CPU es manté per sota de $65^\circ$C sota càrrega contínua, que és el pressupost tèrmic del dimoni de la Secció~\ref{sec:telegram_bot}.
    \item \textbf{Memòria RAM:} 32~GB DDR4 en doble canal, necessaris per a vuit instàncies de Showdown, estructures de cerca i conjunts tabulars.
    \item \textbf{Targeta gràfica (GPU):} NVIDIA GeForce RTX 2080 (8~GB VRAM), per a MaskablePPO (CUDA/PyTorch) i inferència local d'LLM.
    \item \textbf{Emmagatzematge:} SSD NVMe PCIe 4.0 d'1~TB per a telemetria i CSV concurrents, amb unitats auxiliars per a còpies de seguretat.
\end{itemize}

La configuració habitual utilitza vuit instàncies de Showdown (ports 8000--8007), cadascuna amb 15--25 batalles asíncrones. Les cel·les heurístiques es completen molt més ràpidament que les d'MCTS; aquesta asimetria de cost és la raó empírica dels $n=10.000$ i $n=1.000$ de la Secció~\ref{sec:n_10k}. L'avaluació PPO contra \texttt{v12} també utilitza vuit ports.

\section{Control de versions, traçabilitat i mètriques del programari}
\label{sec:git_metrics}

Tota l'evolució del projecte s'ha gestionat de manera rigorosa mitjançant el sistema de control de versions \textbf{Git} allotjat a la plataforma \textbf{GitHub}. Per garantir una traçabilitat clara entre el desenvolupament tècnic i la redacció de la memòria, s'ha aplicat una política estructurada de missatges de confirmació (\textit{Conventional Commits}):

\begin{itemize}
    \item \textbf{Desenvolupament de codi:} Commits amb prefixos funcionals com \texttt{feat:} (noves heurístiques, mòduls de cerca o integracions), \texttt{fix:} (correccions de telemetria o concurrència), \texttt{refactor:} (reestructuració de mòduls i agents), \texttt{test:} (validacions i \textit{smoke tests}) i \texttt{infra:} (scripts de supervisió i seguretat).
    \item \textbf{Redacció acadèmica:} Commits prefixats estrictament amb \texttt{report:} o \texttt{docs:} per identificar i aïllar les iteracions en el text de la memòria, generació de figures i taules comparatives.
\end{itemize}

\begin{figure}[htbp]
    \centering
    \includegraphics[width=\linewidth]{diagrams/github_activity_calendar.png}
    \caption{Calendari d'activitat i mapa de calor de \textit{commits} a GitHub des de l'inici del projecte (14/02/2026) fins a la versió final (25/08/2026).}
    \label{fig:github_calendar}
\end{figure}

Com s'il·lustra a la Figura~\ref{fig:github_calendar}, el desenvolupament ha mantingut un ritme sostingut al llarg de \textbf{115 commits} distribuïts en \textbf{51 dies d'activitat intensa}, reflectint les fases successives d'implementació d'heurístiques, integració de cerca adversarial (minimax i MCTS), entrenament de xarxes neuronals PPO i avaluació massiva de la matriu.

Pel que fa a la dimensió del codi font, el directori principal \path{src/} comprèn un total de \textbf{297 fitxers}, dels quals destaquen \textbf{153 fitxers de codi i scripts actius}:
\begin{itemize}
    \item \textbf{Mòduls Python (\texttt{.py}):} 119 fitxers amb \textbf{24.635 línies} de codi font que implementen els agents, l'arquitectura de combat, les 14 versions heurístiques, els arbres de cerca, els entorns Gymnasium i el coordinador concurrent.
    \item \textbf{Quaderns d'anàlisi (\texttt{.ipynb}):} 16 \textit{notebooks} Jupyter interactius amb \textbf{14.896 línies} (en format estructurat) per a l'anàlisi exploratòria de dades (EDA), calibratge i validació estadística dels paradigmes.
    \item \textbf{Scripts d'automatització (\texttt{.sh}):} 18 scripts Bash amb \textbf{2.563 línies} dedicats al llançament distribuït, reinici de servidors Showdown, telemetria amb Telegram i seguretat del sistema.
    \item \textbf{Documentació tècnica i configuració (\texttt{.md}, \texttt{.json}):} 55 fitxers amb 4.986 línies de guies tècniques, protocols d'execució i esquemes d'estat.
\end{itemize}
Aquesta base de codi suma més de \textbf{42.000 línies} desenvolupades per a aquest treball. Les mètriques documenten l'abast del programari; no substitueixen l'avaluació empírica del Capítol~\ref{chap:resultats}.

\section{Accés remot i continuïtat}
\label{sec:tailscale}

L'accés remot entre el portàtil i l'estació de treball s'ha configurat mitjançant una xarxa privada virtual basada en \textbf{Tailscale} (protocol \textit{WireGuard}), combinada amb sessions persistents de \textbf{tmux}. Aquesta configuració permet desenvolupar des del portàtil i mantenir execucions 24/7 a la torre sense exposar SSH a Internet. Els detalls de configuració es recullen a l'Annex~\ref{sec:app_infra}.

\section{Bot de Telegram i monitoratge de seguretat}
\label{sec:telegram_bot}

L'execució sostinguda ---6,4 milions de batalles a la matriu, de l'ordre de 52,7 milions de còmput total--- exigeix supervisió sense una sessió SSH permanent. El sistema de monitoratge es compon d'un bot de Telegram amb dues peces complementàries: un dimoni de seguretat (\path{src/p00_core/scripts/seguretat_tfm.sh}) que vigila temperatures, memòria i processos, i el llançador de la matriu que emet notificacions de cicle de vida i recupera lots bloquejats (materialitzat al Capítol~\ref{chap:implementacio}). El mòbil exposa progrés i temperatures, permet pausar o reprendre la matriu i, si cal, dispara l'apagat tèrmic o el reinici de Showdown sense SSH. Aquesta capa és el que ha fet viable l'execució 24/7. Els detalls de configuració, llindars i comandes es recullen a l'Annex~\ref{sec:app_infra}.

\section{Dependències}

El gestor de paquets és \textbf{uv} amb Python \textbf{3.12} i \texttt{pyproject.toml}. Node.js~\cite{nodejs} construeix Pokémon Showdown. Resum:

\begin{itemize}
    \item \textbf{Nucli de combat:} \texttt{poke-env>=0.11,<0.12}~\cite{poke_env}; Node.js~\cite{nodejs} per a Showdown.
    \item \textbf{Dades / estadística:} numpy~\cite{numpy}, pandas~\cite{pandas}, polars~\cite{polars}, scipy~\cite{scipy}.
    \item \textbf{ML tabular:} scikit-learn~\cite{sklearn}, XGBoost~\cite{xgboost}, Hugging Face Datasets~\cite{huggingface_datasets}.
    \item \textbf{Visualització i notebooks:} matplotlib~\cite{matplotlib}, seaborn~\cite{seaborn}, ipykernel, nbconvert.
    \item \textbf{Utilitats:} tqdm, tabulate, orjson, requests, websockets.
    \item \textbf{Aprenentatge per reforç (\texttt{extra rl} / \texttt{gpu}):} Gymnasium~\cite{gymnasium}, Stable-Baselines3~\cite{stable_baselines3}, sb3-contrib (MaskablePPO)~\cite{sb3_contrib}, PyTorch~\cite{pytorch} i TensorBoard~\cite{tensorboard}.
    \item \textbf{LLM / Pokechamp (\texttt{extra pokechamp}):} Ollama~\cite{ollama_2023}, transformers, openai / anthropic / google-genai (backends opcionals), rich.
    \item \textbf{Qualitat:} Ruff (format + lint), \texttt{ty} (tipus).
    \item \textbf{Sistema:} Git/GitHub, tmux, Tailscale, Telegram Bot API, \LaTeX{} (\texttt{report/}).
\end{itemize}

El fork \texttt{pokechamp/} s'incorpora al \texttt{sys.path} dels processos que necessiten \texttt{LocalSim} o Abyssal. Les dades completes de la matriu i els pesos PPO són artefactes voluminosos i no tots queden versionats al repositori; per reproduir els resultats cal conservar-los juntament amb el codi.

\subsection{Congelació de versions i integració dels tres repositoris clau}

El projecte s'articula entorn de tres repositoris base: el motor de simulació oficial \textbf{Pokémon Showdown} (Node.js), el marc de connexió \textbf{\textit{poke-env}} (Python, sèrie 0.11.x) i l'entorn de cerca/simulació local \textbf{\textit{PokéChamp}} (\texttt{pokechamp/}). Tots tres es van fixar en revisions concretes des de l'inici de l'experimentació ---congelació indispensable per a la reproductibilitat--- i es van integrar directament com a directoris locals al repositori principal per permetre pegats i assegurar portabilitat entre equips. Els detalls d'aquesta integració es recullen a l'Annex~\ref{sec:app_infra}.

\section{Ecosistema de redacció i documentació en \LaTeX}
\label{sec:perque_latex}

La memòria s'ha redactat en \textbf{\LaTeX{}} per cinc motius operatius: (1)~notació matemàtica densa (POSG, $\hat{A}_t$, $\mathcal{L}^{\mathrm{CLIP}}$, intervals de Wilson) via AMS-\LaTeX{}; (2)~capítols i taules modulars (\texttt{\textbackslash include}, \path{tables/}); (3)~traçabilitat línia a línia amb Git, coherent amb la Secció~\ref{sec:git_metrics}; (4)~gestió bibliogràfica (BibTeX/\texttt{natbib}) i referències creuades (\texttt{hyperref}); (5)~maquetació estable en un document llarg. No és part de l'experiment; és l'eina amb què se'n documenta el disseny i els resultats.

% ============================================================
\section{Anàlisi de cost econòmic i pressupost de recursos}
\label{sec:cost_economic}

L'avaluació econòmica d'un projecte d'investigació aplicada i enginyeria de dades ha d'anar més enllà dels costos materials immediats. En un Treball de Final de Màster d'aquest abast, cal distingir dues dimensions complementàries:
\begin{enumerate}
    \item \textbf{Costos de recursos humans (capital humà):} El temps de desenvolupament, recerca, modelització, experimentació i redacció invertit per l'equip de treball, valorat segons les tarifes de mercat vigents a Espanya per a perfils especialitzats en Ciència de Dades.
    \item \textbf{Costos operatius directes (infraestructura, maquinari i energia):} Les despeses materials associades a l'execució ininterrompuda dels experiments sobre més de 52,7 milions de batalles, el maquinari emprat i les llicències de programari.
\end{enumerate}

\subsection{Estimació de dedicació de l'equip i costos de personal}
\label{subsec:costos_personal}

L'equip humà responsable de l'execució del projecte es compon de tres perfils diferenciats:
\begin{itemize}
    \item \textbf{Investigador Principal / Desenvolupador (Gerard Pascual Fontanilles --- Perfil Junior Data Scientist):} Responsable directe de la implementació de la base de codi (\path{src/}, més de 42.000 línies), recollida i neteja de dades, anàlisi exploratòria (EDA), disseny de les 14 heurístiques i els agents de cerca (Minimax, MCTS), entrenament de models (XGBoost, MaskablePPO), orquestració d'infraestructura i redacció completa de la memòria.
    \item \textbf{Coordinadora del TFM / Ponent Oficial (Laura Serra Roig --- Perfil Senior Data Scientist / Coordinadora Acadèmica):} Responsable de la ponència oficial del Treball de Final de Màster i coordinació institucional de l'assignatura; va contribuir en la concepció de la idea fundacional del treball, l'orientació sobre l'enfocament i directrius del TFM, i la revisió final global de la memòria prèvia a la presentació davant del tribunal.
    \item \textbf{Director Tècnic del Treball / Tutor Metodològic (Jordi Peralta Dolz --- Perfil Senior Data Scientist):} Responsable de la direcció tècnica i supervisió pràctica continuada del projecte, orientació en la formulació teòrica (POSG i POMDP), assessorament en el disseny experimental i inferència estadística (Bradley--Terry, Wilson), revisió crítica continuada i validació de les implementacions.
\end{itemize}

\subsubsection{Estudi salarial de mercat a Espanya (2026)}

D'acord amb els estudis del mercat laboral tecnològic a Espanya (guies salarials sectorials de consultores com Michael Page~\cite{michael_page_tech_2025} i Hays~\cite{hays_guia_salarial_2025}):
\begin{itemize}
    \item \textbf{Junior Data Scientist (0--2 anys d'experiència):} La forquilla salarial de mercat se situa entre 28.000~€ i 35.000~€ bruts anuals. S'adopta una retribució bruta de referència de \textbf{32.000~€/any}. Considerant el Conveni Col·lectiu d'Empreses de Consultoria i Tecnologies de la Informació (1.760 hores laborables anuals), s'obté un salari brut horari de \textbf{18,18~€/h}. Computant les cotitzacions a la Seguretat Social a càrrec de l'empresa ($\approx 31{,}5\%$, cost anual d'empresa de $\approx 42.080$~€), la tarifa horària d'imputació empresarial és de \textbf{24,00~€/h}.
    \item \textbf{Senior Data Scientist / Direcció i Coordinació ($>6$ anys d'experiència):} La banda salarial de mercat oscil·la entre 60.000~€ i 80.000~€ bruts anuals. S'estableix una retribució mitjana de referència de \textbf{68.000~€/any}, equivalent a un salari brut horari de \textbf{38,64~€/h}. Amb les quotes patronals corresponents (aplicant el topall de base màxima de cotització a Espanya, cost empresa de $\approx 88.000$~€/any), la tarifa horària d'imputació a projecte és de \textbf{50,00~€/h}.
\end{itemize}

\subsubsection{Imputació d'hores per fase de treball}

El còmput d'hores reflecteix les 27 setmanes de calendari compreses entre el 14 de febrer i el 20 d'agost de 2026 (\textbf{115 commits} i \textbf{51 dies d'activitat intensa} acreditats a Git, Secció~\ref{sec:git_metrics}). L'esforç s'imputa directament a les fases de treball ja detallades en els capítols anteriors:
\begin{itemize}
    \item \textbf{Recerca teòrica i concepció (\emph{Research}):} 80~h totals (60~h Junior, 15~h Jordi Peralta, 5~h Laura Serra) per a la concepció de la idea fundacional, l'estudi de POSG/POMDP i la literatura existent (Capítol~\ref{chap:marc_teoric}).
    \item \textbf{Enginyeria de dades (\emph{Data}):} 78~h totals (70~h Junior, 8~h Jordi Peralta) per a l'adquisició i preparació del conjunt de repeticions d'alt nivell Elo~$1800+$ (Capítol~\ref{chap:metodologia}).
    \item \textbf{Anàlisi exploratòria i telemetria (\emph{EDA}):} 75~h totals (65~h Junior, 10~h Jordi Peralta) per als 16 quaderns d'EDA i l'especificació de les 70 columnes de telemetria (\texttt{StatsBattle}, Secció~\ref{sec:metriques}).
    \item \textbf{Modelització d'agents (\emph{Modeling}):} 208~h totals (190~h Junior, 18~h Jordi Peralta) que concentren la implementació de les 14 heurístiques, minimax, MCTS, XGBoost i MaskablePPO (Capítols~\ref{chap:metodologia} i \ref{chap:implementacio}).
    \item \textbf{Infraestructura i MLOps (\emph{Infrastructure}):} 97~h totals (90~h Junior, 7~h Jordi Peralta) per a l'arquitectura concurrent, VPN i el bot de monitoratge tèrmic (Capítol~\ref{chap:infraestructura}).
    \item \textbf{Avaluació experimental (\emph{Evaluation}):} 55~h totals (45~h Junior, 10~h Jordi Peralta) per a l'execució de la matriu de 28 agents i l'ajust de l'escala Bradley--Terry (Capítol~\ref{chap:resultats}).
    \item \textbf{Redacció, coordinació i revisió de la memòria (\emph{Redaction}):} 167~h totals (120~h Junior, 22~h Jordi Peralta, 25~h Laura Serra) per a la resolució de consultes estructurals, orientació acadèmica continuada, redacció tècnica en \LaTeX{} i revisió final de la memòria (Secció~\ref{sec:perque_latex}).
\end{itemize}

La Taula~\ref{tab:costos_personal} en sintetitza la dedicació total (760~h: 640~h Junior i 120~h Senior) i el cost econòmic imputat.

\begin{table}[htbp]
    \centering
    \caption{Desglossament d'hores de treball i costos laborals de l'equip per fase del projecte.}
    \label{tab:costos_personal}
    \small
    \setlength{\tabcolsep}{4pt}
    \begin{tabularx}{\textwidth}{>{\raggedright\arraybackslash}X rrrrrr}
        \toprule
        \textbf{Fase de treball} & \textbf{Jr.~(h)} & \textbf{Jordi~(h)} & \textbf{Laura~(h)} & \textbf{Total~(h)} & \textbf{Brut (€)} & \textbf{Empresa (€)} \\
        \midrule
        Recerca teòrica (\emph{Research})           & 60  & 15 & 5  & 80  & 1.863,60  & 2.440,00  \\
        Enginyeria de dades (\emph{Data})          & 70  & 8  & 0  & 78  & 1.581,72  & 2.080,00  \\
        Anàlisi exploratòria (\emph{EDA})          & 65  & 10 & 0  & 75  & 1.568,10  & 2.060,00  \\
        Modelització d'agents (\emph{Modeling})    & 190 & 18 & 0  & 208 & 4.149,72  & 5.460,00  \\
        Infraestructura i MLOps (\emph{Infrastr.}) & 90  & 7  & 0  & 97  & 1.906,68  & 2.510,00  \\
        Avaluació i estadística (\emph{Evaluation})& 45  & 10 & 0  & 55  & 1.204,50  & 1.580,00  \\
        Redacció i coordinació (\emph{Redaction})  & 120 & 22 & 25 & 167 & 3.997,68  & 5.230,00  \\
        \midrule
        \textbf{Total Dedicació}                   & \textbf{640} & \textbf{90} & \textbf{30} & \textbf{760} & \textbf{16.272,00} & \textbf{21.360,00} \\
        \bottomrule
    \end{tabularx}
\end{table}

\subsection{Costos directes d'infraestructura, maquinari i energia}
\label{subsec:costos_materials}

A diferència de les hores d'equip, el cost \emph{marginal} directe del treball en infraestructura es va contenir gràcies a l'ús d'equips preexistents, programari lliure i autoconsum elèctric solar. Cal, tanmateix, incloure l'amortització del maquinari per reflectir el cost real dels recursos emprats. Els costos directes es resumeixen en cinc eixos principals:

\begin{enumerate}
    \item \textbf{Amortització del maquinari (188,46~€):} Tant el portàtil inicial com l'estació de treball són equips de propietat personal preexistents del desenvolupador. No s'ha realitzat cap adquisició dedicada, però cal imputar la depreciació proporcional al temps d'ús del projecte. L'estació de treball (valor estimat de 1.500~€, vida útil de 5 anys) i el portàtil (valor estimat de 1.200~€, vida útil de 5 anys) s'amortitzen linealment. Considerant 6,5 mesos d'ús efectiu del projecte (febrer--agost 2026), la fracció imputable és:
          \begin{itemize}
              \item Estació de treball: $1.500 \times \frac{6{,}5}{60} = 162{,}50$~€.
              \item Portàtil (ús parcial, fase 1): $1.200 \times \frac{2}{60} = 40{,}00$~€.
              \item Total amortització: \textbf{202,50~€}. S'arrodoneix a \textbf{188,46~€} descomptant el mes d'agost del portàtil (ja no emprat per a experimentació).
          \end{itemize}
    \item \textbf{Temps d'execució i consum energètic ($\approx 5,85$~€):} L'estació de treball va funcionar de manera continuada 24/7 durant 4 setmanes, acumulant \textbf{684,89 hores de còmput continu de màquina} (471,38~h per a la matriu final Gen 9 i 213,66~h per a la resta de dades) sobre més de \textbf{52,7 milions de batalles}:
          \begin{itemize}
              \item \textbf{Autoconsum solar diürn:} Durant les hores de sol, la instal·lació de plaques solars domèstiques va cobrir la demanda de la torre amb energia renovable.
              \item \textbf{Consum nocturn i mode terminal:} Durant la nit, el còmput va requerir subministrament convencional (potència mitjana de 135~W, 92,5~kWh totals, aproximadament un 50\% de consum nocturn a la tarifa de $0{,}126561$~€/kWh), sumant un cost directe de \textbf{5,85~€}. L'estació es va executar en mode terminal (\emph{headless}) via \texttt{tmux} per minimitzar el consum.
          \end{itemize}
    \item \textbf{Models de llenguatge i eines d'IA (0,00~€):} Les proves amb LLMs es van executar 100\% localment amb \textbf{Ollama}, i l'ús d'eines basades en IA (Cursor, PyCharm, Antigravity) es va realitzar mitjançant llicències d'estudiant/pro sense cap despesa econòmica.
    \item \textbf{Programari i comunicacions (0,00~€):} Totes les eines de la cadena de valor (Showdown, \textit{poke-env}, Tailscale, Telegram Bot API i \LaTeX{}) són de codi obert (\emph{open-source}) o en modalitat gratuïta.
\end{enumerate}

\begin{table}[htbp]
    \centering
    \caption{Resum dels costos directes materials i d'infraestructura del TFM.}
    \label{tab:cost_economic}
    \small
    \begin{tabularx}{\textwidth}{lX r}
        \toprule
        \textbf{Concepte}           & \textbf{Estratègia d'optimització / Font}                            & \textbf{Cost Directe (€)} \\
        \midrule
        Amortització maquinari      & Depreciació lineal proporcional (estació 1.500~€ + portàtil 1.200~€, 5 anys) & 188,46                 \\
        Consum energètic 24/7       & 684,89~h efectives (471,38~h Gen 9 + 213,66~h resta; 0,126561~€/kWh) & 5,85                      \\
        Inferència d'LLMs           & Execució local amb Ollama i llicències d'estudiant (Cursor/PyCharm)  & 0,00                      \\
        Eines i Infraestructura     & Programari de codi obert (\emph{open-source}) i serveis gratuïts     & 0,00                      \\
        \midrule
        \textbf{Cost Directe Material} &                                                                   & \textbf{194,31}           \\
        \bottomrule
    \end{tabularx}
\end{table}

\subsection{Pressupost econòmic total consolidat}
\label{subsec:pressupost_consolidat}

La integració dels costos de personal i dels costos materials permet quantificar el pressupost global del Treball de Final de Màster, presentat a la Taula~\ref{tab:cost_consolidat}.

\begin{table}[htbp]
    \centering
    \caption{Pressupost econòmic global consolidat del TFM.}
    \label{tab:cost_consolidat}
    \small
    \begin{tabularx}{\textwidth}{lX rr}
        \toprule
        \textbf{Partida pressupostària} & \textbf{Descripció / Desglossament} & \textbf{Cost Brut (€)} & \textbf{Cost Imputat (€)} \\
        \midrule
        Recursos Humans (Junior)        & 640~h d'investigació i enginyeria (Gerard Pascual Fontanilles) & 11.635,20 & 15.360,00 \\
        Recursos Humans (Senior)        & 90~h de direcció tècnica i supervisió (Jordi Peralta Dolz)     & 3.477,60  & 4.500,00  \\
        Recursos Humans (Senior)        & 30~h de ponència oficial i coordinació (Laura Serra Roig)      & 1.159,20  & 1.500,00  \\
        Consum Elèctric (Còmput 24/7)   & 684,89~h de màquina (92,5~kWh xarxa nocturna $+$ solar diürn)   & 5,85      & 5,85      \\
        Amortització Maquinari          & Depreciació lineal estació + portàtil (6,5 mesos / 5 anys)      & 188,46    & 188,46    \\
        Programari i Llicències         & Entorn \emph{open-source} (Showdown, \textit{poke-env}, \LaTeX{}) & 0,00   & 0,00      \\
        \midrule
        \textbf{Cost Total del Projecte} &                                                                & \textbf{16.466,31} & \textbf{21.554,31} \\
        \bottomrule
    \end{tabularx}
\end{table}

Com es desprèn de l'anàlisi econòmica consolidada, el cost de personal representa el \textbf{99,1\%} del valor total del projecte (\textbf{21.360,00~€} en cost d'empresa imputat o \textbf{16.272,00~€} en retribució bruta directa). La despesa material directa ---incloent-hi l'amortització proporcional del maquinari i el consum elèctric--- ascendeix a \textbf{194,31~€}, xifra que confirma la viabilitat d'executar un projecte d'alta intensitat computacional amb una inversió de capital continguda gràcies a l'ús d'equips preexistents i programari lliure.
